\scanpage{391}\mspage{305}

\uuline{Lemme} 4.13.13. Soient $\pi_0,\pi>\pi_0$ des cardinaux infinis, \inserted{tels que $\pi^{\pi_0}=\pi$,} et $\deleted{\unsure{I}}J$ une \inserted{petite} catégorie, grde devant $\pi$. Soit $\deleted{\unsure{J}}I$ l'ens. des sous-catégories $J'$ de $J$ \deleted{telles que} qui sont grdes devant $\pi_0$, et de cardinal $\leq\pi$. Alors $I$ est $\pi_0$-$\pi$-adaptée.

Il reste à voir que pour \unsure{tt} ens. ordonné $A$ grd devant $\pi_0$ et de cardinal $\leq\pi$, si $\unsure{i|r}\mapsto J_r$ est une famille croissante d'éléments de $I$ indexée par $A$, alors

\[
J\overset{\text{\unsure{déf|def}}}{=}\bigcup_{r\in A}J_r=\underset{r\in A}{\underset{\longrightarrow}{\lim}}\,J_r
\]

est dans $I$, i.e. est une sous-catégorie \inserted{de $I$} de cardinal $\leq\pi$ (ce qui est clair) et \uline{grde devant} $\pi_0$. Cela va résulter du

\uline{Lemme} 4.13.13.1 -- Soit $(J_r)_{r\in A}$ un système inductif de catégories grdes devant $\pi_0$, indexé par un ens. ordonné $A$ grd devant $\pi_0$. Alors $J=\underset{\unsure{r\in A}}{\underset{\longrightarrow}{\lim}}\,J_r$ est grd devant $\pi_0$.

% archive sheet number: 390
